\documentclass[11pt,a4paper]{article}

% ===== Paquetes básicos =====
\usepackage[spanish]{babel}
\usepackage[utf8]{inputenc}
\usepackage[T1]{fontenc}
\usepackage{lmodern}
\usepackage{url}
\usepackage{hyperref}
\usepackage{geometry}
\usepackage{setspace}
\usepackage{listings}
\usepackage{xcolor}

\geometry{margin=2.5cm}
\setstretch{1.1}

% Configuración para código
%\lstset{
%	basicstyle=\ttfamily\small,
%	breaklines=true,
%	showstringspaces=false,
%	columns=fullflexible
%}


\lstdefinelanguage{JavaScript}{
	% Palabras clave básicas
	morekeywords=[1]{
		break, case, catch, class, const, continue, debugger, default, delete, do,
		else, export, extends, finally, for, function, if, import, in, instanceof,
		let, new, return, super, switch, this, throw, try, typeof, var, void, while,
		with, yield
	},
	% Literales y tipos
	morekeywords=[2]{
		false, null, true, undefined,
		boolean, number, string, symbol, bigint
	},
	% Objetos y funciones globales
	morekeywords=[3]{
		Array, Boolean, Date, Error, EvalError, Function, JSON, Map, Math,
		Number, Object, Promise, Proxy, RangeError, ReferenceError, RegExp, Set,
		String, SyntaxError, TypeError, URIError, WeakMap, WeakSet,
		console, window, document
	},
	sensitive=true,                % distingue mayúsculas
	% Comentarios
	morecomment=[l]{//},
	morecomment=[s]{/*}{*/},
	morecomment=[s]{/**}{*/},      % estilo JSDoc
	% Cadenas
	morestring=[b]',
	morestring=[b]",
	morestring=[b]`                % template strings
}[keywords, comments, strings]



\lstdefinestyle{jsstyle}{
	language=JavaScript,
	basicstyle=\ttfamily\small,
	keywordstyle=[1]\color{blue}\bfseries,
	keywordstyle=[2]\color{orange!80!black},
	keywordstyle=[3]\color{teal!70!black},
	stringstyle=\color{green!50!black},
	commentstyle=\color{gray},
	numberstyle=\tiny\color{gray},
	numbers=left,
	stepnumber=1,
	breaklines=true,
	showstringspaces=false,
	columns=fullflexible
}


\lstdefinestyle{terminalstyle}{
	basicstyle=\ttfamily\small,
	backgroundcolor=\color{black!5},
	keywordstyle=\color{black},
	commentstyle=\color{black},
	stringstyle=\color{black},
	showstringspaces=false,
	breaklines=true,
}



\title{Detectando tipos de integración en un caso sencillo}
\author{Juan Alvarado}
\date{}

\begin{document}
	
	\maketitle
	
	\section*{Intrucciones}
	
	%5.1 Manual de laboratorio de la Práctica 1
	%“Detectando tipos de integración en un caso sencillo”
	
	\begin{itemize}
		\item Formar 4 equipos y elegir un contexto organizacional concreto (por ejemplo, biblioteca universitaria, clínica pequeña, tienda en línea de productos digitales).
	 	
		\item 	Identificar al menos seis sistemas o aplicaciones que podrían existir en ese contexto (por ejemplo, sistema de inventario, punto de venta, CRM, portal web, sistema de facturación, app móvil).
		 	
		\item 	Dibujar un diagrama simple con cajas para cada sistema y flechas que representen intercambios de información conocidos o razonables.
	 	
		\item 	Para cada flecha, describir en una frase qué se intercambia (por ejemplo, “ventas diarias”, “datos de clientes”, “citas médicas confirmadas”).
		 	
		\item 	Clasificar cada intercambio como principalmente de datos, funcional o semántico; si hay duda, anotar el motivo de la duda para discutirlo.
		 	
		\item 	Elegir tres integraciones que parezcan críticas y escribir brevemente qué podría salir mal si la integración fallara en cada uno de los tres niveles (datos inconsistentes, funciones que no se coordinan, significados distintos).
		 	
		\item 	Preparar una mini‑presentación de 5 minutos por equipo donde muestren su diagrama y discutan al menos un caso de cada tipo de integración.
	\end{itemize}
	
	
\end{document}
